\documentclass{subfiles}
\begin{document}
    As for Huffman encoding, the arithmetic encoding procedure just discussed
        has an issue: the probabilities needs to be known before the encoding begins.
        A question then arises: can we do better as for the adaptive Huffman encoding?
        That is, can we make arithmetic encoding adaptive? 
        In the present section we discuss how we can achieve adaptability.

    The idea is the same as that of adaptive Huffman, but instead of 
        updating the tree, we update the known probabilities.
        Algorithmically speaking, it begins by assuming that all symbol have the same probability;
        that is, the \(n\) subintervals have the same width.
        As for the static case, once a symbol is read it selects the corresponding subinterval,
        and increases that symbol probability. The selected 
        interval is then divided in \(n\) subintervals with width proportional to the probabilities.

    \begin{example*}
        Let \(S = abcb\) be the text to encode. 
            Before any character is read, the probabilities are: \(p_{a} = p_{b}
                = p_{c} = \sfrac{1}{3}\).
        Once the first character (\(a\)) is read, we select the range \([0, \sfrac{1}{3}]\)
        and update the probabilities as \(p_{a} = \sfrac{1}{2}, p_{b} = p_{c} = \sfrac{1}{4}\).
        We repeat this process until the end of the text, 
        encoding \(S\) within the range \([\sfrac{43}{180}, \sfrac{11}{45}]\).
        \subfile{../../TikZ/Figure - Example of adaptive arithmetic encoding.tex}        
    \end{example*}

    The decoding step is symmetrical: instead of updating the probabilities once a 
        symbol is read, these are updated each time a new one is decoded.
        Apart from this caveat, both the encoding and decoding procedures 
        are equal to \autoref{proc:arithmeticEncoding} and \autoref{proc:arithmeticDecoding}.    
\end{document}
